\begin{pastexam}[label={px:gw3-104-ds-05}]{關務三等資訊處理 104 年 第五題}
（申論；15 分）

學生的學號格式是（N$_{1}$N$_{2}$N$_{3}$N$_{4}$N$_{5}$N$_{6}$N$_{7}$），假設儲存空間為 99，請用數字分析法（Digital Analysis），
分別以學號為鍵值（Key）雜湊（Hashing）出其資料儲存的位址。數字的分布曲度（Skewness）設為 $sk$，則
$sk_{i} = \sum\limits_{j=0\sim 9} |\,a_{ij} - 1\,|$，其 $a_{ij}$ 表示 N$_{i}$ 出現的個數。
\begin{itemize}[leftmargin=2.2em, nosep]
\item[㈠] 請依下列五位學生的學號算出其 $sk_{i}$ 值。（10 分）

\begin{tabular}{@{}l}
\texttt{Student 1 ID: 0392018} \\
\texttt{Student 2 ID: 0392124} \\
\texttt{Student 3 ID: 0392238} \\
\texttt{Student 4 ID: 0252714} \\
\texttt{Student 5 ID: 0392468}
\end{tabular}
\item[㈡] 請寫出此五位學生儲存的位址。（5 分）
\end{itemize}
\end{pastexam}

\begin{pastsoln}
答案：㈠ $sk_1 = 13$、$sk_2 = 11$、$sk_3 = 11$、$sk_4 = 13$、$sk_5 = 5$、$sk_6 = 7$、$sk_7 = 11$。 ㈡ 取 N$_5$N$_6$：
學生 1--5 的位址依序是 $01$、$12$、$23$、$71$、$46$。

\textbf{解讀。} 題目的「$a_{ij}$ 表示 N$_i$ 出現的個數」意思是：在這五個學號中，第 $i$ 位數字是 $j$ 的有幾個。
$\sum_j a_{ij} = 5$。 若某一位的數字分布得很平均（$0 \sim 9$ 各出現約一次），每個 $|a_{ij} - 1|$ 都小，$sk_i$ 就小；
若五個學號這一位都一樣，$sk_i$ 最大。

\textbf{數字分析法}（digit analysis）也是一種雜湊函數（資結書第~12.1~節介紹了除法、中間平方、摺疊等做法）：事先知道全部的鍵時，逐位檢查數字的
分布，丟掉分布偏斜（$sk$ 大）的位數，留下分布最平均的幾位組成位址。 儲存空間 $99$ 格（位址 $00 \sim 98$），需要兩位數，
所以取 $sk$ 最小的兩位。

\textbf{㈠ 逐位計算。}
\begin{center}
\begin{tabular}{cll}
\toprule
位 & 五個學號在這一位的數字 & $sk_i = \sum_{j=0}^{9} |a_{ij} - 1|$ \\
\midrule
N$_1$ & 0, 0, 0, 0, 0 & $|5 - 1| + 9 \times |0 - 1| = 4 + 9 = 13$ \\
N$_2$ & 3, 3, 3, 2, 3 & $|4 - 1| + |1 - 1| + 8 \times 1 = 3 + 0 + 8 = 11$ \\
N$_3$ & 9, 9, 9, 5, 9 & $|4 - 1| + |1 - 1| + 8 \times 1 = 11$ \\
N$_4$ & 2, 2, 2, 2, 2 & $4 + 9 = 13$ \\
N$_5$ & 0, 1, 2, 7, 4 & 五個數字各一次：$5 \times 0 + 5 \times 1 = 5$ \\
N$_6$ & 1, 2, 3, 1, 6 & $1$ 出現兩次：$|2 - 1| + 3 \times 0 + 6 \times 1 = 7$ \\
N$_7$ & 8, 4, 8, 4, 8 & $|3 - 1| + |2 - 1| + 8 \times 1 = 2 + 1 + 8 = 11$ \\
\bottomrule
\end{tabular}
\end{center}

\textbf{㈡ 位址。} $sk$ 最小的兩位是 N$_5$（$5$）和 N$_6$（$7$），取 N$_5$N$_6$ 當位址：
\begin{center}
\begin{tabular}{lcc}
\toprule
學生 & 學號 & 位址（N$_5$N$_6$） \\
\midrule
Student 1 & 0392\textbf{01}8 & $01$ \\
Student 2 & 0392\textbf{12}4 & $12$ \\
Student 3 & 0392\textbf{23}8 & $23$ \\
Student 4 & 0252\textbf{71}4 & $71$ \\
Student 5 & 0392\textbf{46}8 & $46$ \\
\bottomrule
\end{tabular}
\end{center}
五個位址都不同，沒有碰撞，而且都在 $00 \sim 98$ 之內。 （若兩位數取到 $99$，超出 $99$ 格的範圍，可再 $\bmod 99$。）
\end{pastsoln}
